\chapter{Dead Code and Control Dependence}
\label{ch:ch15-dead-code-control-dependence}

Optimizers create garbage as a side effect of doing their job.
Constant folding leaves dead definitions behind, branch folding leaves dead edges behind, and inlining leaves dead parameter setups behind.
Each corpse costs code size, instruction cache, and analyst attention in later passes.
Worse, dead code can keep live code pessimistic: a dead store blocks a load motion, a dead branch doubles a join, and a dead call blocks alias refinement.
The problem of this chapter is to delete code that cannot affect observable behavior, without deleting code that can.
We combine backward liveness with control reasoning so that data-dead and control-dead code both disappear.
Section~\ref{sec:ch15-dead} defines deadness, Section~\ref{sec:ch15-control} adds control dependence, and Section~\ref{sec:ch15-worked} deletes by hand on Figure~\ref{fig:ch15-cfg}.

\section{What dead means}
\label{sec:ch15-dead}

An instruction is \term{dead} when removing it cannot change any observable behavior of any execution.
A more operational test uses liveness: a pure definition is dead when no live use reads it on any path.
A store is dead when no load that may read it is live, accounting for aliasing conservatively.
A branch is dead when its condition is constant, in which case the untaken edge and its exclusive blocks may be unreachable.
A call is dead only when it provably lacks side effects and its result is unused, which few calls satisfy.
The \enquote{provably} qualifier carries the whole safety burden, and every deletion in this chapter traces back to it.
Table~\ref{tab:ch15-patterns} lists the deletion patterns with the proof each one demands.

Consider \code{t = a + b} followed by \code{use a} with no later mention of \code{t}.
The addition defines a temporary nobody reads, performs no side effect, and raises no trap under the type guards.
Deleting it preserves every output, fault, and divergence the program had.
Now change the operation to \code{t = x / y} where \code{y} might be zero.
Deleting the division removes a potential trap, turning a failing execution into a succeeding one.
That changes observable behavior, so the division is live for trapping purposes even when its value is unused.
Beginners delete the second case because their tests never divide by zero, and the miscompile surfaces only on hostile input.
The C standard classifies several such operations as trapping or undefined, so portable deletion must consult it, see~\cite{wg14-2011-c11}.

\begin{irlisting}
; Join with a dead definition, a live use, and a foldable branch
; B1 branches on flag, B2 defines dead t and live u, B3 defines u
B1: br flag, B2, B3
B2: t = a + b
B2: u = a + 1
B2: jmp B4
B3: u = a + 2
B3: jmp B4
B4: use u
\end{irlisting}

The listing holds one clearly dead instruction and one temptation.
The definition of \code{t} in $B2$ is dead because no path uses \code{t} afterward and integer addition has no effect beyond its value.
The definitions of \code{u} are both live because $B4$ uses \code{u} on every path.
Deleting \code{t} is safe today and profitable forever, since later passes stop analyzing it.
Deleting either definition of \code{u} would corrupt the use at $B4$ on some path.
A liveness pass over this graph proves the first claim and refutes the second, as Section~\ref{sec:ch15-worked} shows numerically.
Algorithm~\ref{alg:ch15-adce} then generalizes the hand reasoning to loops and control dependence.
\section{Control dependence and open questions}
\label{sec:ch15-control}

Data liveness alone cannot delete a branch, because branches define no values.
A block is \term{control dependent} on a branch when the branch decides whether the block executes.
Formally the dependence uses post-dominators: $Y$ depends on $X$ when $X$ has two successors whose paths to exit rejoin differently around $Y$.
Computing it needs the \term{post-dominator} tree over the reversed graph, where entry and exit swap roles.
The \term{control dependence graph} then links each branch to the blocks it governs, and dead-edge removal consults it before deleting.
A branch whose condition is constant governs nothing worth keeping on the dead arm, so its control dependents go dark together.
Figure~\ref{fig:ch15-cfg} shows $B2$ and $B3$ both dependent on $B1$, which is why folding $B1$ can delete either arm wholesale.

Unreachable code needs this machinery rather than plain liveness.
An infinite loop with no live values still governs everything after it by never falling through.
A branch on a constant still selects one arm even when both arms compute live values locally.
Marking-based deletion therefore starts from truly live roots, such as returns, stores to escaped memory, and effectful calls, then marks backward through data and control edges together.
Anything unmarked after the sweep is dead, including whole regions guarded by constant branches.
The sweep must treat volatile operations, atomic fences, and input-output calls as roots unconditionally.
It must also treat potential traps conservatively unless a prior proof discharges them.
Safonov treats this root discipline as part of the trusted-deletion argument, see~\cite{safonov-2010-trustworthy}.

\begin{algorithm}
\caption{Conservative marking-based dead-code elimination}\label{alg:ch15-adce}
\begin{pseudocode}
\Require CFG $G$ with explicit effects and traps, and complete reaching-definition links
\Ensure unused pure definitions removed; all original control flow retained
\State $Live \gets$ roots: all terminators, observable uses, stores, and operations not proved pure, total, and nontrapping
\State $MarkedB \gets$ blocks holding roots
\State $W \gets$ all root instructions
\While{$W \neq \emptyset$}
\State remove some $i$ from $W$
\ForAll{operands $v$ used by $i$}
\ForAll{definitions $d$ of $v$ that can reach $i$}
\If{$d$ unmarked}
\State mark $d$ and add $d$ to $W$
\EndIf
\EndFor
\EndFor
\ForAll{branches $c$ that $i$ is control dependent on}
\If{$c$ unmarked}
\State mark $c$ and add $c$ to $W$
\EndIf
\EndFor
\EndWhile
\State delete unmarked pure definitions; retain all terminators and blocks
\end{pseudocode}
\end{algorithm}

Algorithm~\ref{alg:ch15-adce} marks live instructions rather than deleting dead ones directly.
It seeds roots that observable behavior requires, then walks backward through data uses and control conditions together.
A definition becomes live when a live instruction reads it, and a branch becomes live when a live instruction depends on it.
Termination follows because each instruction marks at most once and each edge is examined on marking.
Correctness follows because every deleted instruction is pure, unread, and ungoverning by construction of the unmarked set.
This conservative sweep retains every terminator, so it preserves loops and termination behavior without relying on an exit-reaching post-dominator model.
Branch deletion is a separate extension requiring a proof about control flow, phi selection, effects, traps, and divergence.
Morgan presents deletion passes in this mark-then-sweep shape inside the optimization pipeline, see~\cite{morgan-1998-optimizing}.
\section{Worked example with marking}
\label{sec:ch15-worked}

We run Algorithm~\ref{alg:ch15-adce} by hand on Figure~\ref{fig:ch15-cfg} with root \code{use u} in $B4$.
With no successor uses, $\LiveOut{B4}=\emptyset$; the observable use inside $B4$ makes $\LiveIn{B4}=\{u\}$.
Before deletion, ordinary liveness at $B2$ includes both $a$ and $b$ because the dead assignment \code{t = a + b} still reads them.
After deleting that pure assignment, only $a$ remains live on entry to $B2$.
Block $B3$ defines \code{u} from \code{a}, giving live-in $\{a\}$.
Block $B1$ uses \code{flag} for the branch and passes $\{a, flag\}$ upward.
Table~\ref{tab:ch15-patterns} classifies each instruction from these sets, and the marking trace below confirms the classification constructively.

\begin{figure}
\begin{tikzpicture}
\node[basic block] (b1) {B1: br flag};
\node[basic block, below left=1.1cm and 0.7cm of b1] (b2) {B2: t dead, u live};
\node[basic block, below right=1.1cm and 0.7cm of b1] (b3) {B3: u = a + 2};
\node[basic block, below=2.4cm of b1] (b4) {B4: use u};
\draw[cfg edge] (b1) -- (b2);
\draw[cfg edge] (b1) -- (b3);
\draw[cfg edge] (b2) -- (b4);
\draw[cfg edge] (b3) -- (b4);
\end{tikzpicture}
\caption{The dead-code diamond where \code{t} in $B2$ is dead but both definitions of \code{u} stay live through the use in $B4$.}
\label{fig:ch15-cfg}
\end{figure}

\begin{table}
\caption{Deletion patterns with the safety proof each one requires before removal.}
\label{tab:ch15-patterns}
\begin{tabular}{lcc}
\toprule
Pattern & Safety proof needed & Action \\
\midrule
Pure def, no live use & liveness fixed point & delete definition \\
Dead branch arm & constant condition & fold to jump \\
Unreachable block & no live control path & delete block \\
Effectful or trapping op & no pureness proof & keep operation \\
\bottomrule
\end{tabular}
\end{table}

Marking seeds \code{use u} in $B4$ as live and enqueues it.
It marks both definitions of \code{u} in $B2$ and $B3$ through data edges, since each reaches the use along some executable path.
It marks the branch in $B1$ through control edges, since the live $B2$ and $B3$ definitions both depend on it.
It marks \code{a} and \code{flag} transitively as operands of marked instructions.
It never marks \code{t = a + b} because no marked instruction reads \code{t} and \code{t} governs nothing.
The sweep deletes exactly that addition and keeps everything else.
The \enquote{one deletion} outcome is the point: aggressive marking removes what liveness refutes and nothing more.
A second scenario folds $B1$ on constant true first, after which $B3$ loses live control and both its instructions go dark in the next sweep.
The cascade shows why deletion and folding iterate together rather than running once each.

The same example warns against deleting the tempting \code{u} copy in $B3$.
Its value differs from the $B2$ copy, and the join in $B4$ reads whichever arm executed.
Removing it would leave the right path reading an undefined version, which renaming would expose as a missing phi argument.
Control dependence explains the error structurally: $B3$ is governed by $B1$, $B4$ is not governed by either arm, and the use in $B4$ needs both governors live.
Data reasoning and control reasoning agree, which is the invariant to check whenever they seem to disagree.
When they do disagree, the implementation should keep the code and file a precision bug rather than delete and hope.
The exercises ask you to construct exactly such a disagreement with aliasing.

\section{Implementation notes and cost}
\label{sec:ch15-implementation}

The conservative algorithm needs reaching-definition links and a sound side-effect oracle; because all terminators are roots, its control-dependence loop adds no new marks and may be omitted.
Aggressive branch deletion needs additional control-dependence and termination reasoning.
Given $I$ instructions and $D$ explicit dependency links, marking and sweeping cost $O(I+D)$ time and storage, excluding construction of those links.
Non-SSA reaching-definition links can be numerous; neither their construction nor post-dominator computation is made linear merely by using a worklist.
Iteration across folding cascades repeats sweeps, but each sweep deletes at least one instruction or edge, bounding the number of sweeps.
Production compilers run a cheap trivial pass first, for single-use chains and constant branches, and reserve aggressive marking for post-folding cleanup.
Appel shows where each layer sits relative to register allocation and code emission, see~\cite{appel-1998-modern}.

The oracle is the trust boundary, not the graph walk.
A pure annotation on a call must be verified against the callee body or a trusted specification, never against the call name alone.
An alias query that declares a store dead must cover every load on every path, including exceptional and concurrent paths the CFG omits.
A trap-deletion decision must cite the type and range proof that discharges the fault, not the absence of failures in testing.
These are the rules that keep deletion from becoming miscompilation with good benchmark scores.
The C language definition supplies the trapping, sequencing, and volatility facts the oracle consults, see~\cite{wg14-2011-c11}.
Cytron and colleagues give a primary treatment connecting control-dependence graphs with dominance-frontier machinery~\cite{cytron-1991-ssa}.

\section{Exercises}
\label{sec:ch15-exercises}

\begin{enumerate}
\item \label{ex:ch15-hand-mark} Run Algorithm~\ref{alg:ch15-adce} on Figure~\ref{fig:ch15-cfg} after changing $B2$ to \code{t = x / y}.
Assume the division may trap and consult Table~\ref{tab:ch15-patterns} for the trapping row.
Show the new root set, the marking order, and whether the division survives.
Then add a range proof that \code{y} is nonzero and rerun, showing which oracle answer flips the outcome.
\item \label{ex:ch15-alias-trap} Add \code{store u to p} in $B2$ and \code{load q} in $B4$ with unknown aliasing between \code{p} and \code{q}.
Decide by hand whether the store is dead under a may-alias oracle and under a must-alias oracle.
State which oracle a safe deleter must use and what test would expose the aggressive choice as wrong.
Relate your answer to the volatility rules in the C definition.
\item \label{ex:ch15-cascade-measure} Implement marking plus branch folding for diamonds chained $N$ deep, with $N$ set to $10$, $100$, and $1000$.
Seed a constant at the first branch so each fold deadens one arm and exposes the next constant.
Count sweeps to quiescence and instructions deleted per sweep for the joint iteration versus one marking pass followed by one folding pass.
Report the counts and recommend the production layering.
\end{enumerate}

\section{Further reading}
\label{sec:ch15-reading}

Appel locates dead-code elimination among the passes that must run before allocation, see~\cite{appel-1998-modern}.
Morgan treats control-sensitive deletion as part of the loop and redundancy pipeline, see~\cite{morgan-1998-optimizing}.
Safonov explains the root and oracle discipline that keeps deletion trustworthy, see~\cite{safonov-2010-trustworthy}.
The C language definition fixes the observable behavior that deletion must preserve, see~\cite{wg14-2011-c11}.
